Recent work on chain-of-thought (CoT) faithfulness reports single aggregate numbers (e.g., DeepSeek-R1 acknowledges hints 39\% of the time~\cite{chen2025reasoning}), implying that faithfulness is an objective, measurable property of a model. This paper demonstrates that it is not. Three classifiers---a regex-only detector, a two-stage regex-plus-LLM pipeline, and an independent Claude Sonnet~4 judge---are applied to 10,276 influenced reasoning traces from 12 open-weight models spanning 9 families and 7B to 1T parameters. On identical data, these classifiers produce overall faithfulness rates of 74.4\%, 82.6\%, and 69.7\%, respectively, with non-overlapping 95\% confidence intervals. Per-model gaps range from 2.6 to 30.6 percentage points; all are statistically significant (McNemar's test, $p < 0.001$). The disagreements are systematic, not random: inter-classifier agreement measured by Cohen's $\kappa$ ranges from 0.06 (``slight'') for sycophancy hints to 0.42 (``moderate'') for grader hints, and the asymmetry is stark---for sycophancy, 883 cases are classified as faithful by the pipeline but unfaithful by the Sonnet judge, while only 2 go the other direction. Classifier choice can also reverse model rankings: Qwen3.5-27B ranks 1st under the pipeline but 7th under the Sonnet judge; OLMo-3.1-32B moves in the opposite direction, from 9th to 3rd. The root cause is that different classifiers operationalize the same criterion (``load-bearing'' hint acknowledgment) as different constructs---lexical \emph{mention} versus epistemic \emph{dependence}---and these constructs yield divergent measurements on the same behavior. These results demonstrate that published faithfulness numbers cannot be meaningfully compared across studies that use different classifiers, and that future evaluations should report sensitivity ranges across multiple classification methodologies rather than single point estimates.
